%!TEX root = ../thesis.tex
% ******************************* Thesis Appendix A ****************************
\chapter{Where Does Value Come From?} 

\ifpdf
    \graphicspath{{Appendix1.2/Figs/Raster/}{Appendix1.2/Figs/PDF/}{Appendix1.2/Figs/}}
\else
    \graphicspath{{Appendix1.2/Figs/Vector/}{Appendix1.2/Figs/}}
\fi

We open this thesis with the idea that utility is to choices as reproductive fitness is to innate behaviour. As genes give rise to phenotypes to maximise their reproductive success, some process in the brain is generating some internal sense of value that guides the choices an organism makes and the sense of reward, and thus reinforcement, they accrue from it\footnote{This framing that most, if not all, of conscious human behaviour is making economic choices is partly done for hyperbole and may be contentious, but is not in itself wrong. ‘Economic imperialism’ has regularly claimed all purposeful (human) behaviour as under the field of economics (see \cite{stiglerEconomicsImperialScience1984}). If all behaviour is economic, one might wonder why so little is written on the ‘psychological/ethological imperialism’ of the minds that perform these choices. A neuroscientific imperialism of the brain that must surely create those minds is outlined in \cite{bujoldAdaptiveEconomicsNeuroethological2019}}. In the introduction to the thesis, we introduce the idea of utility, and the formalisation of Expected Utility Theory by von Neumann and Morgenstern (\cite{vonNeumann1944}).

Although we think the results presented in this thesis have merit and show that the BDM mechanism can be used to glimpse the physiological substrate of utility, we are clearly only at the beginning of an arduous course of research to fully uncover the computations performed by organisms to make economic decisions; currently, we only have a glimpse of these processes.

We use this appendix to place ideas and discussions that took place over the course of this project with other members of the Schultz laboratory. These boil down mainly to the idea that if dopamine is 'a neural substrate for prediction and reward' (Schultz et al., 1997) then what does that mean. Obviously, these questions are not limited to dopaminergic neurons (see Appendix C), and indeed, at the end of this Appendix, we discuss a similarly formulated question that focuses on the prefrontal cortices; 'Where does value come from?’ (\cite{juechemsWhereDoesValue2019})\footnote{Throughout the main text of the thesis, we try to only refer to ‘utility’ as the quantity maximised by our monkey’s making economic choices. Here we break from that somewhat and talk about 'value' to fit the literature we will be discussing, and also because we want to use this appendix to digress and discuss the ideal of ‘value’ rather than observed measurements of ‘utility’}.

\section{Where Does Value Come From?}

For most people, considerations of value likely begin and end with the price of various consumer goods. Prices and the willingness of consumers to pay them are the signal by which the majority of the world's economy is arranged; however, to a consumer, they can seem as if they are divinely ordained. Some higher power has decided that good $X$ costs £$Y$ (give or take some small error between shops and over time). In order to explain how economies work, much of the history of economic thought has sought to explain this key signal: What leads an object to be labelled with a specific price?

‘Classical’ economic approaches to this question initially settled on the idea that value was intrinsic to the good being sold. Adam Smith originally formulated the idea that the cost of labour put into offering a good for exchange is the \textit{'ultimate standard by which the value of commodities… can be estimated and compared'}, later termed the ‘Labour Theory of Value’ (\cite{smith1776inquiry}).  This was expounded upon by Ricardo, who relaxed the theory, stating that the price of a good is merely proportional to labour (\cite{stiglerRicardo93Labor1958}). This idea of the value of a good being embodied by the cumulative labour necessary to produce it is a fundamental property of Marxist economic thought, with the difference between the value (i.e., the labour supplied) and the price of a good representing the surplus value (commonly termed S), that is, the exploitation of labour (\cite{marx1867capital}).

At the very beginning of this thesis, we skip over this century of economic thought between Bernoulli’s formulation of utility versus value and its rise to prominence in ‘neoclassical’ economics, among which is the idea that value is generated internally by a consumer. Furthermore, we only briefly touch on neoclassical economic thought to discuss Edgeworth’s invention of indifference curves for choices (\cite{edgeworth1881mathematical}) and do not consider the prices of goods in the market.

Two other proponents of this school were William Stanley Jevons, who in 1862 (published \cite{jevonsMathematicalTheoryPolitical1874}) outlined the idea of marginal utility as the driving force behind prices, and Léon Walras, who took this idea further with his idea of a ‘general equilibrium’ (\cite{walras1954elements}). In such an equilibrium, each subject maximises their utility while ‘clearing’ a market (i.e. there is no leftover supply or demand in the economic environment). Such a toy equilibrium is shown for two monkeys trading bananas and apples in Figure \ref{fig:indifference}, extending Edgeworth's representation of indifference curves illustrated in Figure \ref{fig:indifference}.

\begin{figure}[htbp!] 
\centering    
\includegraphics[width=1\textwidth]{AppendixB/Figs/walras.png}
\caption[Neoclassical theories of utility provide the framework for Walrasian equilibria.]{\textbf{Neoclassical theories of utility provide the framework for Walrasian equilibria.} Imagine our monkey from \ref{fig:indifference} choosing between apples and bananas has a partner, and we name these monkeys ‘Monkey U’ and ‘Monkey V’ (in line with the results chapters of this thesis). Both monkeys have indifference curves between apples and bananas, as in Figure \ref{fig:indifference} (blue and green; in this case symmetrical, though in practice probably unique to each organism). An initial allocation of apples and bananas ‘X’ leaves Monkey V with more bananas than apples, and vice versa for Monkey U. These allocations lie upon indifference curves of low utility for both animals (dashed lines). A Walrasian equilibrium can be found where the monkeys trade the good they have a surplus of for the other, up to a point which lies upon the indifference curve representing the maximum achievable joint utility for both monkeys (‘E’; solid lines). At point E, both curves intersect along a tangent (red, dotted line), the slope of which is equal to the ratio of the price for bananas and the price for apples (in terms of each other) for both monkeys.}
\label{fig:walras}
\end{figure}

The marginal utility of each good, represented by the indifference point, is the key contributor to the demand of each monkey for the two fruits provided which, along with the initial allocation (the supply), will lead to prices. This role of marginal utility in determining the demand for goods is the key to solving Adam Smith’s ‘diamond paradox’. Although the first unit of water supplied to a dehydrated person is worth more than the first unit of diamonds, in general, water is plentiful, so an additional unit of diamonds contains more marginal utility.

The capstone of Walras’ theory of value were extensions by Ken Arrow and Gérard Debreu in their 1954 model of general equilibria in markets, for which the former won the Nobel Prize in Economics in 1972 (\cite{arrowExistenceEquilibriumCompetitive1954}). Arrow and Debreu showed that, with some assumptions, a competitive market should arrive at prices which balance the supply and demand (i.e. the utility) of goods across the entire range of products.

In general, the neoclassical school of economics and their ‘marginal revolution’ showed that goods do not have intrinsic value. As Carl Menger stated (\cite{menger1950principles}):

\begin{quote}
Value is thus nothing inherent in goods, no property of them, nor an independent thing existing by itself. It is a judgment economizing men make about the importance of the goods  at their disposal for the maintenance of their lives and well-being. Hence, value does not exist outside the consciousness of men\footnote{Though we prefer a simpler statement misquoted from a Clint Black song:
“[value] isn’t something [an object] has; it’s something that we do”.}
\end{quote}

If value is only found in consciousness, it seems clear to us that to find some true, objective value of something, one would need to begin by looking into the brains of humans and other animals.

\section{Do We Like Value or Do We Want It?}

The idea of utility as something internal to an individual’s mind remains a key tenet of economic thought. However, it raises the question of where this utility arises from. Throughout the main body of this thesis, we dodge this question, regarding animals as 'utility-maximising' computers that will seek the greatest reward (defined as anything that reinforces the behaviour to receive said reward). For the study of whether and how monkeys perform auction tasks and the correlates of dopaminergic neuron firing during these tasks, this is a satisfactory definition. For this reason, neuroeconomic experiments use objective reward dimensions, in this case, the magnitude of the juice offered, but also, commonly, the probability or even the quantified nutritional content (\cite{grabenhorstRepresentationOralFat2014}). More certain and larger rewards containing necessary nutrients are dominantly rewarding over those that are less likely to be delivered, are smaller, or do less to satisfy basic physiological needs. We 'want' to maximise the rewards we obtain, in the same way that a computer program can 'want' a higher score on a video game (\cite{mnihPlayingAtariDeep2013}) or to win a game of chess (\cite{silverGeneralReinforcementLearning2018}).

However, in practice, humans and other animals do not spend their entire lives chasing rewards. When possible, that is, when the organism has accumulated enough of these base desires, organisms can expend temporal, social, or even monetary capital. In fact, in the colloquial usage of the word ‘reward’, many of the experiences that people list as the most rewarding (such as travelling and time spent with partners and family) revolve around experiencing pleasure without regard to physiological needs. Even when these ‘wanting’ and ‘liking’ overlap in a single good, it is relatively clear that they are not the same thing. For example, when writing, the author often treats himself to a large bag of sugary sweets. Though they are able to space the first few out over many minutes and enjoy the pleasant sweet taste of them, invariably by the end of the first hour or so this is replaced by a craving to continue eating even as the 'liking' has been diminished or perhaps even reversed.

This idea of ‘wanting’ and ‘liking’ being dissociable concepts first appeared in neuroscience to explain the lack of impact of 6-OHDA dopaminergic neurons on the ability of rats to taste (\cite{berridgeTasteReactivityAnalysis1989})\footnote{For a more recent review of the idea and criticism that goes beyond this appendix, see \cite{berridge2013neuroeconomics}.}. Although the lesions led to deficits in the drinking and eating behaviour of rats, they were still able to differentiate between tastes and retained their palate. That is, rats still 'liked' the taste of pleasurable food sources, but no longer ‘wanted’ them (in contrast to the gluttonous author and their sweets). Similar findings have been reported in humans, where catecholamine treatments (\cite{leytonCocaineCravingEuphoria2005}) and pimozide application (\cite{brauerHighDosePimozide1997}) reduces the craving for drugs such as cocaine and ecstasy, but not the euphoric feeling elicited when the drugs are taken.

This raises a computational question: if dopamine does not itself carry pleasure, what role does it play in the reward circuitry we describe elsewhere in this thesis? Early attempts to reconcile 'wanting' with prediction-error accounts treated incentive salience as a form of cached, Pavlovian-learned value (\cite{mcclureComputationalSubstrateIncentive2003}; \cite{zhangNeuralComputationalModel2009}), while later work has argued more directly that dopamine's incentive signal is a distinct computation from a value-prediction error, shaped by the animal's momentary physiological state as well as by what has been learned (\cite{berridgePredictionErrorIncentive2012}).

An intrinsic problem of trying to measure this experienced utility (even more so than for our simple definition of utility used in the main body of this text!) is that, by its nature, these are highly subjective feelings. A bag of sweets can be rationally said to have more utility than a single sweet to the author as he enters a cinema foyer, but exactly how pleasurable is the act of chewing each? fMRI studies have attempted to study this by presenting a stimulus and then modulating the experience of it. For example, by 'revealing' that an odour comes from expensive cheese or an unwashed body (\cite{dearaujoCognitiveModulationOlfactory2005}), or that a glass of wine is from a cheaper or more expensive vintage (\cite{plassmannWhatCanAdvertisers2007}). In such studies, changes in the subject's experience of stimuli are correlated with BOLD activity in the orbitofrontal cortex, suggesting a role for this in the 'liking' component of utility.

Although the current evidence for these distinct streams of utility presents intriguing possibilities, more research is needed to establish more definitive causal links between experimental manipulations and the dissociation of 'liking' versus 'wanting' behaviours in laboratory animals. We acknowledge the inherent challenges in such investigations (how can one obtain an accurate behavioural report of an animal's momentary pleasure?) and find merit in the conceptual distinction between predicted and experienced utility. For instance, observations of human Parkinson's patients treated with deep brain stimulation to recover dopaminergic activity suggest that self-stimulation addresses a \textit{pre hoc} craving rather than providing a pleasurable \textit{post hoc} experience (\cite{berridgeBuildingNeurosciencePleasure2011}). This framework offers promising avenues for future neuroeconomic research into the multifaceted nature of utility and the processes in the brain that produce these ideas.

More recent work continues to develop this picture. Berridge has argued that incentive salience can separate sharply from an animal's own predictions of outcome value, even producing desire for outcomes remembered to be unpleasant (\cite{berridgeSeparatingDesirePrediction2023}), and a thirty-year retrospective on the underlying incentive-sensitization theory finds the core wanting/liking dissociation well supported (\cite{robinsonIncentiveSensitizationTheoryAddiction2025}). At the level of individual cells, medium spiny neurons in the nucleus accumbens have recently been shown to carry dissociable signals, with one population tracking stimulus salience and another tracking prediction error more conventionally (\cite{zachryD1D2Medium2024}), although how uniformly dopamine itself contributes to valence coding across the striatum remains debated (\cite{rodriguesDopamineNucleusAccumbens2025}; \cite{guillauminDisentanglingRoleNAc2023}).

\section{Neuroethology of Value}

If wanting and liking are separable, do we value goods because we like them or because we want them? A recent perspective by Keno Juechems and Chris Summerfield tends toward the latter, namely that the utility of a good or choice to an organism reflects the distance it closes between the current state of the animal (hunger, thirst, boredom, etc.) and some multidimensional goal state (\cite{juechemsWhereDoesValue2019}). Such models are ‘homeostatic’ in that they incentivise organisms to choose options as if taking the shortest distance back to an ideal set point where basic needs are met (\cite{keramatiHomeostaticReinforcementLearning2014}; \cite{oreillyGoalDrivenCognitionBrain2014}).

The authors argue that their model is supported by recent data in which participants must balance the accumulation of assets along two distinct dimensions (filling a zoo with virtual lions and elephants; \cite{juechemsNetworkComputingValue2019}). The ventromedial prefrontal cortex\footnote{This paper develops out of a human fMRI paper which focus on the cortex. See \cite{bogaczDopamineRoleLearning2020} for an interpretation of dopaminergic neurons in this context.} tracked progress toward redressing the balance of animals rather than the amount of goods (animals) received per se.

When studying human economic behaviour, we have access to a well-studied and characterised signal of value: the price of a good. In this thesis, we spend a whole chapter defining the BDM auction task partly to explain and justify its use, but to make it clear that we can experimentally access this continuous signal. Therefore, this project only deals with one (type of) animal and one manipulation (the volume of the good offered) of one kind of good (liquid rewards); whether this is a useful way to study value is questionable. The early marginal revolutionists realised that these prices arise as a function of the fungibility of goods, for example, with our hypothetical monkeys trading fruits in Figure \ref{fig:walras}.

\cite{haydenCaseEconomicValues2021} take a similar approach when arguing against the case for an encoding of economic values neither in the orbitofrontal cortex nor anywhere else in the brain. The authors argue that most naturalistic decisions faced by animals are between goods that are more unlike than they are alike. Instead, the authors argue for a 'heuristic' model of decision-making. The authors cite work by \cite{slovicConstructionPreference1995} which shows that preferences are constructed at the point at which they are interrogated, rather than as some lookup comparing the values of the options presented. We would not go as far as Hayden and Niv; the recordings made in Chapter 4 of this thesis (Figures \ref{fig:fractal_display} and \ref{fig:fractal_display_regression}) seem to act 'as if' dopaminergic midbrain cells are encoding the value of an offered good, if only on an artificially constructed task.

These more naturalistic approaches to neuroeconomic questions also raise the question of why and how the brain is able to make economic decisions. The theories of utility we have described across the first chapter of this thesis and its appendices are beautiful pieces of mathematical logic; the brain is a wet mess of organic tissue. Over the first decades of the field, neuroeconomists have tended to use this framework as a 'top-down' approach to search for where value signals might exist in the brain (\cite{wuEmergingStandardNeurobiological2018}) instead of 'bottom-up' approaches, even as failures of expected utility theory to animal behaviour are well publicised, e.g. Prospect Theory's subjective weighting of probabilities in both humans \cite{kahnemanProspectTheoryAnalysis1979} and monkeys \cite{bujoldComparingUtilityFunctions2022}. Interestingly, recent work has attempted to take such an ethological framework when studying economic decision-making at both the behavioural (\cite{glimcherEfficientlyIrrationalIlluminating2022}) and neuronal \cite{polaniaEfficientCodingSubjective2019} level.

In summary, we see 3 open dichotomies in the study of what value is and where it comes from:

\begin{itemize}
    \item \textbf{wanting vs. liking} and whether specific regions of the brain or neuronal populations encode either. 
    \item \textbf{goal-directed vs. utility maximising decision-making} and the extent to which either are relevant to naturalistic economic decision-making.
    \item \textbf{top-down vs. bottom-up models of economic decision-making} and how neuroscience can inform economic models of utility rather than become slaves to them.
\end{itemize}

As with many biological 'dichotomies', we take the banal opinion that the productive way forward is to take both sides of these debates seriously; the study of why we value things is a vast question and we are not nearly done. Indeed, with respect to the final point, one recent model of 'Expected Subjective Value Theory' (ESVT) seeks to marry expected utility theory with biological plausibility (\cite{glimcherExpectedSubjectiveValue2023}).